
\section{The Hermite integration method} \label{ch:hermite}


Each particle is completely specified by its mass $m$,
position $\mathbf r_0$, and velocity $\mathbf v_0$,
where the sub\-script $0$ denotes an initial value at a
time $t_0$.
The equation of motion for a particle $i$ is given by
its momentary acceleration $\mathbf a_{0,i}$ due to
all other particles and its time derivative
$\dot{\mathbf a}_{0,i}$ as
%
\begin{eqnarray}
\mathbf a_{0,i} &=& -\sum_{i \neq j} G m_j\frac{\mathbf R}{R^3} ,
                 \label{eq:accel0}\\
%
\dot{\mathbf a}_{0,i} &=& -\sum_{i \neq j} G m_j
          \left [ \frac{\mathbf V}{R^3}
                  + \frac{3\mathbf R(\mathbf V \cdot \mathbf R)}{R^5}
          \right ] , \label{eq:accel1}
\end{eqnarray}
%
where $G$ is the gravitational constant;
$\mathbf R = \mathbf r_{0,i} - \mathbf r_{0,j}$ is
the relative coordinate;
$R = |\mathbf r_{0,i} - \mathbf r_{0,j}|$ the modulus;
and $\mathbf V = \mathbf v_{0,i} - \mathbf v_{0,j}$
the relative space velocity to the particle $j$.

The Hermite scheme employed in NBODY6++ follows the trajectory
of the particle by firstly ``predicting'' a new position
and new velocity for the next time step $t$.
A Taylor series for $\mathbf r_i(t)$ and $\mathbf v_i(t)$
is formed:
%
\begin{eqnarray}
\mathbf r_{p,i}(t) &=& \mathbf r_0
                + \mathbf v_0(t-t_0)
                + \mathbf a_{0,i}\frac{(t-t_0)^2}{2}
                + \dot{\mathbf a}_{0,i}\frac{(t-t_0)^3}{6} ,
         \label{eq:predr} \\
%
\mathbf v_{p,i}(t) &=&  \mathbf v_0
                + \mathbf a_{0,i}(t-t_0)
                + \dot{\mathbf a}_{0,i}\frac{(t-t_0)^2}{2} .
          \label{eq:predv}
\end{eqnarray}
%
The predicted values of $\mathbf r_p$ and $\mathbf v_p$,
which result from this simple Taylor series evaluation,
using the force and its time derivative at $t_0$,
do not fulfil the requirements for an accurate
high--order integrator;
they just give a first approximation to $\mathbf r_1$
and $\mathbf v_1$ at the upcoming time $t_1$.
Even if the time step, $t_1 - t_0$, is chosen
impracticably small, a considerable error will
quickly occur, let alone the inadequate
computational effort.
Therefore, an improvement is made by the Hermite
interpolation which approximates the higher
accelerating terms by another Taylor series:
%
\begin{eqnarray}
\mathbf a_i(t) &=& \mathbf a_{0,i}
                   + \dot{\mathbf a}_{0,i} \cdot (t - t_0)
                   + \frac{1}{2}\mathbf a_{0,i}^{(2)} \cdot (t-t_0)^2
                   + \frac{1}{6}\mathbf a_{0,i}^{(3)} \cdot (t-t_0)^3 ,
          \label{eq:accp0} \\
%
\dot{\mathbf a}_i(t) &=& \dot{\mathbf a}_{0,i}
                   + \mathbf a_{0,i}^{(2)} \cdot (t-t_0)
                   + \frac{1}{2}\mathbf a_{0,i}^{(3)} \cdot (t-t_0)^2 .
          \label{eq:accp1}
\end{eqnarray}
%
Here, the values of $\mathbf a_{0,i}$ and $\dot{\mathbf a}_{0,i}$
are already known, but a further derivation of equation
(\ref{eq:accel1}) for the two missing orders on the right
hand side turns out to be quite cumbersome.
Instead, one determines the additional acceleration terms
from the predicted (``provisional'') $\mathbf r_p$ and
$\mathbf v_p$;
we calculate their acceleration and time derivative
according to the equations (\ref{eq:accel0}) and
(\ref{eq:accel1}) anew and call these new terms
$\mathbf a_{p,i}$ and $\dot{\mathbf a}_{p,i}$,
respectively.
Because these values ought to be generated by the former
high--order terms also (which we avoided),
we put them into the left--hand sides of (\ref{eq:accp0})
and (\ref{eq:accp1}).
Solving equation (\ref{eq:accp1}) for $\mathbf a_{0,i}^{(2)}$,
then substituting it into (\ref{eq:accp0}) and simplifying yields
the third derivative:
%
\begin{eqnarray}
\mathbf a_{0,i}^{(3)}
        &=& 12\frac{\mathbf a_{0,i} - \mathbf a_{p,i}}%
                   {(t - t_0)^3}
           + 6\frac{\dot{\mathbf a}_{0,i} + \dot{\mathbf a}_{p,i}}%
                   {(t - t_0)^2} . \label{eq:accel3}
\end{eqnarray}
%
Similarly, substituting (\ref{eq:accel3}) into
(\ref{eq:accp0}) gives the second derivative:
%
\begin{eqnarray}
\mathbf a_{0,i}^{(2)}
        &=&  -6\frac{\mathbf a_{0,i} - \mathbf a_{p,i}}%
                    {(t - t_0)^2}
             -2\frac{2\dot{\mathbf a}_{0,i} + \dot{\mathbf a}_{p,i}}%
                    {t - t_0}.  \label{eq:accel2}
\end{eqnarray}
%
Note, that the desired high--order accelerations are found
just from the combination of the low--order terms for
$\mathbf r_0$ and $\mathbf r_p$.
We never derived higher than the first derivative, but
achieved the higher orders easily through (\ref{eq:accel0})
and (\ref{eq:accel1}).
This is called the Hermite scheme.


Previously, a four--step Adams--Bashforth--Moulton integrator
was used (especially in NBODY5, \cite{a_85}),
however, the new Hermite scheme allows twice as large timesteps
for the same accuracy.
Also its storage requirements are less
\cite{m_91b}, \cite{ma_92}, \cite{a_99a}, \cite{a_99b}.

Finally, we extend the Taylor series for $\mathbf r_i(t)$
and $\mathbf v_i(t)$, eqs.~(\ref{eq:predr}) and
(\ref{eq:predv}), by two more orders,
and find the ``corrected'' position $\mathbf r_{1,i}$
and velocity $\mathbf v_{1,i}$ of the particle $i$
at the computation time $t_1$ as
%
\begin{eqnarray}
\mathbf r_{1,i}(t)
        &=&  \mathbf r_{p,i}(t)
           + \mathbf a_{0,i}^{(2)}\frac{(t - t_0)^4}{24}
           + \mathbf a_{0,i}^{(3)}\frac{(t - t_0)^5}{120} , \\
%
\mathbf v_{1,i}(t)
        &=&  \mathbf v_{p,i}(t)
           + \mathbf a_{0,i}^{(2)}\frac{(t - t_0)^3}{6}
           + \mathbf a_{0,i}^{(3)}\frac{(t - t_0)^4}{24} .
\end{eqnarray}
%
The integration cycle for other upcoming steps may now be
repeated from the beginning, eqs.\ (\ref{eq:accel0}) and
(\ref{eq:accel1}).
The local error in $\mathbf r$ and $\mathbf v$ within the two time
steps $\Delta t = t_1 - t_0$ is expected to be of order
${\cal O}(\Delta t^5)$, the global error for a fixed physical
integration time scales with ${\cal O}(\Delta t^4)$ \cite{m_91a}.

